% (c)~2014 Claudio Carboncini - claudio.carboncini@gmail.com
% (c)~2014 Dimitrios Vrettos - d.vrettos@gmail.com
% (c) 2015 Daniele Zambelli daniele.zambelli@gmail.com

% \section{Esercizi}
% 
% \subsection{Esercizi dei singoli paragrafi}
% 
% \subsubsection*{\numnameref{sec:01_}}
% 
% \begin{esercizio}
% \label{ese:D.19}
% testo esercizio
% \end{esercizio}
% 
% \begin{esercizio}\label{ese:03.1}
% Consegna:
% \begin{enumeratea}
% \item 
% \end{enumeratea}
% \end{esercizio}
% 
% \subsection{Esercizi riepilogativi}
% 
% \begin{esercizio}
% \label{ese:D.19}
% testo esercizio
% \end{esercizio}
% 
% \begin{esercizio}\label{ese:03.1}
% Consegna:
% \begin{enumeratea}
% \item 
% \end{enumeratea}
% \end{esercizio}

\section{Esercizi}
\subsection{Esercizi dei singoli paragrafi}
% \subsection*{4.1 - Risoluzione delle disequazioni di secondo grado}
\numnameref{sec:diseq_secondo_grado}

\begin{esercizio}
\label{ese:4.7}
Rappresentare nel riferimento cartesiano ortogonale le seguenti parabole.

\begin{htmulticols}{2}
\begin{enumeratea}
\item \( y=-3x^2+x \)
\item \( y=\dfrac 1 2x-2x+\dfrac 3 2 \)
\item \( y=x^2+x-1 \)
\item \( y=x^2-x+1 \)
\end{enumeratea}
\end{htmulticols}
\end{esercizio}

\begin{esercizio}
\label{ese:4.8}
Rappresentare nel riferimento cartesiano ortogonale le seguenti parabole.

\begin{htmulticols}{2}
\begin{enumeratea}
\item \( y=-3x^2+3 \)
\item \( y=x^2+4x+3 \)
\item \( y=x^2+\dfrac 3 5 \)
\item \( y=-\dfrac 2 5x^2+4x-\dfrac 1 5 \)
\item \( y=-\dfrac 1 2 x^2-4x-1 \)
\end{enumeratea}
\end{htmulticols}
\end{esercizio}


\affiancati{.39}{.59}{
\begin{esercizio}
\label{ese:4.9}
Per ciascun grafico di parabola \(y=ax^2+bx+c\) indica il segno del primo 
coefficiente e del discriminante, la natura dei suoi zeri (reali distinti, 
reali coincidenti, non reali), il segno della funzione.
\end{esercizio}
}{
\immagine[.7]{Cinque parabole in un piano cartesiano.}{\parabole}
}

\begin{esercizio}[*]
\label{ese:4.1}
Risolvi le seguenti disequazioni di secondo grado.

\begin{htmulticols}{2}
\begin{enumeratea}
\item \(5x^2-6x \leqslant 0\) 
\sol{0 \leqslant x \leqslant \dfrac{6}{5}}
\item \(5x^2>0\) \sol{x\neq 0}
\item \(x^2+x>0\) \sol{x<-1 \sor x>0}
\item \(x^2 \leqslant 0\) \sol{x=0}
\item \(3x^2 \leqslant -1\) \sol{\emptyset }
\item \(x^2-9>0\) \sol{x_1<-3 \sor x>3}
\item \(-x^2+5 \leqslant 0\) 
\sol{x \leqslant -\sqrt 5 \sor x \geqslant \sqrt 5}
\item \(2x^2-3x+1>0\) \sol{x<\dfrac 1 2 \sor x>1}
\item \(-x^2+3x \geqslant 0\) \sol{0 \leqslant x \leqslant 3}
\item \(3x^2+x-2>0\) \sol{x_1<-1 \sor x>\dfrac 2 3}
\item \(x^2-4>0\) \sol{x_1<-2 \sor x>2}
\item \(\dfrac 4 3x^2-\dfrac 1 3x-1<0\) \sol{-\dfrac 3 4<x<1}
\item \(x^2-8 \leqslant 0\) \sol{-2\sqrt 2 \leqslant x \leqslant 2\sqrt 2}
\item \(x^2-4x+9>0\) \sol{\R}
\item \(x^2-6x+8 \leqslant 0\) \sol{2 \leqslant x \leqslant 4}
\item \(x^2+3x-4 \geqslant 0\) \sol{x \leqslant -4 \sor x \geqslant 1}
\item \(x^2-9x+18<0\) \sol{3<x<6}
\item \(-2x^2 \geqslant 0\) \sol{x=0}
\item \(2x^2+5x+4 \leqslant 0\) \sol{\emptyset}
\item \(x^2+1>0\) \sol{\R}
\item \(x^2-8x+15 \geqslant 0\) \sol{x \leqslant 3 \sor x \geqslant 5}
\item \(x^2+5>0\) \sol{\R}
\end{enumeratea}
\end{htmulticols}
\end{esercizio}

\begin{esercizio}[*]
\label{ese:4.4}
Risolvi le seguenti disequazioni di secondo grado.

\begin{htmulticols}{2}
\begin{enumeratea}
\item \(x^2+x \geqslant 0\) \sol{x \leqslant -1 \sor x \geqslant 0}
\item \((x+1)^2 \geqslant 0\) \sol{\R}
\item \(x^2>1\) \sol{x<-1 \sor x>1}
\item \(2x^2-6<0\) \sol{-\sqrt 3<x<\sqrt 3}
\item \(-x^2-1 \leqslant 0\) \sol{\R}
\item \(9-4x^2 \leqslant 0\) 
\sol{x \leqslant -\dfrac{3}{2} \sor x \geqslant +\dfrac{3}{2}} 
\item \(3x-2x^2>0\) 
\sol{x \leqslant -\sqrt{\dfrac{3}{2}} \sor 
   x \geqslant +\sqrt{\dfrac{3}{2}}} 
\item \(x^2 \geqslant 0\) \sol{x = 0}
\item \(2x^2+4>0\) \sol{\R} 
\item \(x^2-x-2>0\) \sol{x \leqslant 1 \sor x \geqslant +2} 
\item \(x^2+11x+30 \leqslant 0\) \sol{x \leqslant -6 \sor x \geqslant -5}
\item \(3x^2+2x \leqslant 0\) 
\sol{-\dfrac{2}{3} \leqslant x \leqslant 0} 
\end{enumeratea}
\end{htmulticols}
\end{esercizio}

% \subsection*{4.2 - Risoluzione grafica di una disequazione di secondo grado}
% 
% \begin{esercizio}
% \label{ese:4.10}
% Risolvere graficamente le seguenti disequazioni di secondo grado.
% \begin{htmulticols}{2}
% \begin{enumeratea}
% \item \( 2x^2+3x-1<0 \)
% \item \( x^2-5x+6 \leqslant 0 \)
% \item \( x^2-3x-4>0 \)
% \item \( x^2-6x+5 \geqslant 0 \)
% \end{enumeratea}
% \end{htmulticols}
% \end{esercizio}
% 
% \begin{esercizio}
% \label{ese:4.11}
% Risolvere graficamente le seguenti disequazioni di secondo grado.
% \begin{htmulticols}{2}
% \begin{enumeratea}
% \item \( 6x^2+x-2>0 \)
% \item \( 15x^2+x-6 \leqslant 0 \)
% \item \( -x^2+1 \geqslant 0 \)
% \item \( x^2-\dfrac 1 4>0 \)
% \end{enumeratea}
% \end{htmulticols}
% \end{esercizio}
% 
% \begin{esercizio}
% \label{ese:4.12}
% Risolvere graficamente le seguenti disequazioni di secondo grado.
% \begin{htmulticols}{2}
% \begin{enumeratea}
% \item \( x^2-\dfrac 1 4x \leqslant 0 \)
% \item \( x^2+2x \leqslant 0 \)
% \item \( x^2+2x+1 \leqslant 0 \)
% \item \( x^2+x+1<0 \)
% \end{enumeratea}
% \end{htmulticols}
% \end{esercizio}
% 
% \begin{esercizio}[*]
% \label{ese:4.13}
% Risolvi le disequazioni di secondo grado con il metodo algebrico o con 
% quello 
% grafico.
% \begin{htmulticols}{3}
% \begin{enumeratea}
% \end{enumeratea}
% \end{htmulticols}
% \end{esercizio}
% 
% \begin{esercizio}[*]
% \label{ese:4.14}
% Risolvi le disequazioni di secondo grado con il metodo algebrico o con 
% quello 
% grafico.
% \begin{htmulticols}{3}
% \begin{enumeratea}
% \item \(-x^2+4x+3>0\)
% \item \(x^2+4x+4<0\)
% \item \(x^2-x+1<0\)
% \item \(x^2-\dfrac 1 9 \geqslant 0\)
% \item \(9x^2+3x-2 \leqslant 0\)
% \item \(2x^2+5<0\)
% \end{enumeratea}
% \end{htmulticols}
% \end{esercizio}
% 
% \begin{esercizio}[*]
% \label{ese:4.15}
% Risolvi le disequazioni di secondo grado con il metodo algebrico o con 
% quello 
% grafico.
% \begin{htmulticols}{2}
% \begin{enumeratea}
% \item \(4x-x^2 \geqslant 0\)
% \item \(9x^2+10x+1 \leqslant 0\)
% \item \(0,01x^2-1>0\)
% \item \(1,\bar 6x^2-2x \leqslant 0\)
% \end{enumeratea}
% \end{htmulticols}
% \end{esercizio}
% 
% \begin{esercizio}[*]
% \label{ese:4.16}
% Risolvi le disequazioni di secondo grado con il metodo algebrico o con 
% quello 
% grafico.
% \begin{htmulticols}{2}
% \begin{enumeratea}
% \item \(\dfrac 1 2x^2-\dfrac 1 8>0\)
% \item \(4x^2+\dfrac 5 3x-1 \leqslant 0\)
% \item \(x^2+x+\sqrt 2>0\)
% \item \(x^2+2\sqrt 2x+2>0\)
% \end{enumeratea}
% \end{htmulticols}
% \end{esercizio}

% \paragraph{4.13.} a)~\(x \leqslant -\dfrac 3 2 \sor x \geqslant \dfrac 3 
% 2\),\quad 
% b)~\(0<x<\dfrac 3 
% 2\),\quad c)~ \(\R\),\quad d)~ \(\R\),\; e)~ \(x<-1 \sor x>2\),\; f)~ 
% \(-6 \leqslant x \leqslant 
% -5\)
% 
% \paragraph{4.14.} a)~\(2-\sqrt 7<x<2+\sqrt 7\),\quad b)~\(\emptyset \) 
% ,\quad 
% c)~\(\emptyset \) ,\quad d)~\(x \leqslant -\dfrac 1 3 \sor x \geqslant 
% \dfrac 1 3\),\quad 
% e)~\(-\dfrac 
% 2 3 \leqslant x \leqslant \dfrac 1 3\),\quad f)~\(\emptyset \)
% 
% \paragraph{4.15.} a)~\(0 \leqslant x \leqslant 4\),\quad b)~ \(-1 \leqslant 
% x<-\dfrac 1 
% 9\),\quad 
% c)~\(x<-10 \sor x>10\),\quad d)~\(0 \leqslant x<\dfrac 6 5\)
% 
% \paragraph{4.16.} a)~\(x<-\dfrac 1 2 \sor x>\dfrac 1 2\),\quad b)~
% \(-\dfrac 3 4 \leqslant x \leqslant 
% \dfrac 1 3\),\quad c)~\(\R\),\quad d)~\(\R-\{\sqrt 2\}\)

\begin{esercizio}[*]
\label{ese:4.17}
Risolvi le disequazioni di secondo grado.
% 
\begin{enumeratees}
\item \(x^2+6x-2>0\) \sol{x<-3-\sqrt{11} \sor x>-3+\sqrt{11}}
\item \(x^2-5x+3 \geqslant 0\)
\sol{x \leqslant \dfrac{5-\sqrt{13}} 2 \sor x \geqslant \dfrac{5+\sqrt{13}} 
2}
\item \(3x^2-\dfrac 2 3x-1 \leqslant 0\) 
\sol{\dfrac{1-2\sqrt 7} 9 \leqslant x \leqslant \dfrac{1+2\sqrt 7} 9}
\item \(x^2-3x-\dfrac 5 2<0\) 
\sol{\dfrac{3-\sqrt{19}} 2<x<\dfrac{3+\sqrt{19}} 2}
\item \(x^2-4x-9 \leqslant 0\) \sol{2-\sqrt{13} \leqslant x \leqslant 
2+\sqrt{13}}
\item \(12x^2-3 \geqslant 4x(2x-1)\) 
\sol{x \leqslant -\dfrac 3 2 \sor x \geqslant \dfrac 1 2}
\item \(2x^2-11x-6 \geqslant 0\) 
\sol{x \leqslant -\dfrac 1 2 \sor x \geqslant 6}
\item \((3x+1)^2>(2x-1)^2\) 
\sol{x<-2 \sor x>0}
\item \((x+1)(x-1)^2>x^3\) 
\sol{-\dfrac{\sqrt 5+1} 2<x<\dfrac{\sqrt 5-1} 2}
\item \((x+3)(x+2)<-(x+2)^2\)
\sol{-\dfrac 5 2<x<-2}
\item \(\dfrac{x+1} 2+\dfrac{(x+1)(x-1)} 4>x^2-1\)
\sol{-1<x<\dfrac 5 3}
\item \((x+1)^3-(x+2)^2>\dfrac{2x^3-1} 2\)
\sol{x<\dfrac{1-\sqrt{21}} 4 \sor x>\dfrac{1+\sqrt{21}} 4}
\item \((x-2)(3-2x) \geqslant x-2\)
\sol{1 \leqslant x \leqslant 2}
\item \((3x+1)\left(\dfrac 5 2+x\right) \leqslant 2x-1\)
\sol{-\dfrac 7 6 \leqslant x \leqslant -1}
\item \(\dfrac{x^2+16} 4+x-1<\dfrac{x-3} 2\)
\sol{\emptyset}
\item \(\dfrac{3x-2} 2<x^2-2\)
\sol{x<-\dfrac 1 2 \sor x>2}
\item \(\dfrac{x-3} 2-\dfrac{x^2+2} 3<1+x\)
\sol{\R}
\end{enumeratees}
\end{esercizio}

\begin{esercizio}[*]
\label{ese:4.20}
Risolvi le disequazioni di secondo grado.

\begin{enumeratea}
\item \((x+4)^2+8 \geqslant \dfrac{x-1} 3\)
\sol{\R}
\item \(\left(\dfrac{x-1} 3-\dfrac x 6\right)^2 \leqslant (x+1)^2\)
\sol{x \leqslant -\dfrac 8 5 \sor x \geqslant -\dfrac 4 7}
\item \(\dfrac 1 2\left(x-\dfrac 2 3\right)^2+x\left(x-\dfrac 2 3\right)
\left(x+\dfrac 2 3\right)>x^3-\dfrac x 2\left(x-\dfrac 2 3\right)-\dfrac 
8{27}\)
\sol{x<\dfrac 2 3 \sor x>\dfrac 7 9}
\item \(3x-5+(1-3x)^2>(x-2)(x+2)\)
\sol{x<0 \sor x>\dfrac 3 8}
\item \(\dfrac{x-2} 3-(3x+3)^2>x\)
\sol{-\dfrac{29}{27}<x<-1}
\item \((x-4)^2+(2-x)^2-2(2x+17)>4(x+5)(3-x)+(x+1)^2\)
\sol{x<-3 \sor x>5}
\item \((x-2)^3-x^3>x^2-4\)
\sol{\dfrac{6-2\sqrt 2} 7<x<\dfrac{6+2\sqrt 2} 7}
\item \((2-x)^3-(2-x)^2<\dfrac{3-4x^3} 4\)
\sol{\emptyset}
\item \((x+2000)^2+x+2000<2\)
\sol{-202<x<-199}
\item \(\dfrac{\left(2x-1\right)^3-8x} 2-\dfrac{\left(2x+1\right)^2-15} 
4 \leqslant 
4x\left(x-1\right)^2-6\)
\sol{}
\item \(\dfrac{(3-x)^2} 2-1 \geqslant -\dfrac{x^2-4} 4\) 
\sol{x \leqslant 2-\dfrac{\sqrt 6} 3 \sor x \geqslant 2+\dfrac{\sqrt 6} 
3}
\item \(\left(\dfrac x 2+1\right)^2-2x>\dfrac 5 4\left(\dfrac x 2-1\right)\)
\sol{}
\item \((x+1)^2>(x-1)^2+(x+2)^2+4x\)
\sol{\emptyset }
\item \(\dfrac{x^2} 4+x<\dfrac{x+3} 4+\dfrac x 2-\dfrac{1-\dfrac x 2} 2\)
\sol{-1<x<1}
\end{enumeratea}
\end{esercizio}

\begin{esercizio}
\label{ese:4.23}
Il monomio \(16x^2\) risulta positivo per:

\fbox{A}\quad \(x>16\)\qquad \fbox{B}\quad \(x>\dfrac 
1{16}\)\qquad\fbox{C}\quad 
\(x<-4 \sor 
x>16\)\qquad\fbox{D}\quad \(x\in \R\)\qquad\fbox{E}\quad \(x\in \R_0\)

\end{esercizio}

\begin{esercizio}
\label{ese:4.24}
Il binomio \(16+x^2\) risulta positivo per:

\fbox{A}\; \(x>-16\)\quad \fbox{B}\; \(-4<x<4\) \quad\fbox{C}\; \(x\in 
\R-\{-4,4\}\) 
\quad\fbox{D}\; \(x\in \R\) \quad\fbox{E}\; \(x<-4 \sor x>4\)
\end{esercizio}

\begin{esercizio}
\label{ese:4.25}
Il binomio \(16-x^2\) risulta positivo per:

\fbox{A}\; \(x>-16\)\quad \fbox{B}\; \(-4<x<4\) \quad\fbox{C}\; \(x\in 
\R-\{-4,4\}\) 
\quad\fbox{D}\; \(x\in \R\) \quad\fbox{E}\; \(x<-4 \sor x>4\)
\end{esercizio}

% \begin{esercizio}
% \label{ese:4.26}
% Spiegate sfruttando il metodo grafico la verità della proposizione: 
% ``nessun 
% valore della variabile \(a\) rende il polinomio \((3+a)^2-(2a+1)\cdot 
% (2a-1)-(a^2+2a+35)\) positivo''.
% \end{esercizio}

% \subsection*{4.3 - Segno del trinomio a coefficienti letterali}
% 
% \begin{esercizio}[*]
% \label{ese:4.27}
% Risolvi e discuti le seguenti disequazioni.
% \begin{htmulticols}{2}
% \begin{enumeratea}
% \item \(x^2-2{kx}+k^2-1>0\)
% \item \(3x^2-5{ax}-2a^2<0\)
% \end{enumeratea}
% \end{htmulticols}
% \end{esercizio}
% 
% \begin{esercizio}[*]
% \label{ese:4.28}
% Risolvi e discuti le seguenti disequazioni.
% \begin{htmulticols}{2}
% \begin{enumeratea}
% \item \(4x^2-4x+1-9m^2<0\)
% \item \(2x^2-3{ax}<0\)
% \end{enumeratea}
% \end{htmulticols}
% \end{esercizio}
% 
% \begin{esercizio}[*]
% \label{ese:4.29}
% Risolvi e discuti le seguenti disequazioni.
% \begin{htmulticols}{2}
% \begin{enumeratea}
% \item \(x^2-2{tx}-8t^2>0\)
% \item \((1-s)x^2+9>0\)
% \end{enumeratea}
% \end{htmulticols}
% \end{esercizio}
% 
% \begin{esercizio}[*]
% \label{ese:4.30}
% Risolvi e discuti le seguenti disequazioni.
% \begin{htmulticols}{2}
% \begin{enumeratea}
% \item \((m-1)x^2-{mx}>0\)
% \item \({kx}^2-(k+1)x-3 \geqslant 0\)
% \end{enumeratea}
% \end{htmulticols}
% \end{esercizio}
% 
% \paragraph{4.27.} a)~\(x<k-1 \sor x>k+1\),\, b)~\(a=0\to \emptyset\) 
% \(a>0\to -\dfrac 1 
% 3a<x<2a\) \(a<0\to 2a<x<-\dfrac 1 3a\)
% 
% \paragraph{4.28.} a)~\(m=0\to \emptyset\) \(m>0\to 
% \dfrac{1-3m}2<x<\dfrac{1+3m} 2\) 
% \(m<0\to \dfrac{1+3m} 2<x<\dfrac{1-3m} 2\),\quad b)~\(a=0\to \emptyset\) 
% \(a>0\to 0<x<\dfrac 3 2a\) \(a<0\to \dfrac 3 2a<x<0\)
% 
% \paragraph{4.29.} a)~\(t=0\to x\neq 0\) \(t>0\to -2t<x<4t\) \(t<0\to 
% 4t<x<-2t\),\quad 
% b)~\(s \leqslant 1\to \R\) \(s>1\to \dfrac{-3}{\sqrt{k-1}}<x<\dfrac 
% 3{\sqrt{k-1}}\)
% 
% \paragraph{4.30.} a)~\(m=0\to \emptyset\) \(m=1\to x<0\) \(0<m<1\to \dfrac 
% m{m-1}<x<0\) 
% \(m<0\to 0<x<\dfrac m{m-1}\) \(m>1\to x<0 \sor x>\dfrac m{m-1}\)

% \begin{esercizio}
% \label{ese:4.31}
% Trovare il segno del trinomio \(t=(1-m)x^2-2{mx}-m+3\) al variare del 
% parametro~\( m \)
% \end{esercizio}

% \newpage %--------------------------------------

% \subsection*{4.4 - Disequazioni polinomiali di grado superiore al secondo}
\numnameref{sec:diseq_scomposizione}

\begin{esercizio}
\label{ese:4.32}
Data la disequazione 
\(\left(x^2-x\right) \cdot \left(2x^2+13x+20\right)<0\) 
verificare che nessun numero naturale appartiene all'insieme soluzione. 
C'è qualche numero intero nell'\( \IS \)? È vero che l'\( \IS \) è 
formato dall'unione di due intervalli aperti di numeri reali?
\end{esercizio}

\begin{esercizio}
\label{ese:4.33}
Dopo aver scomposto in fattori il polinomio \(p(x)=2x^4-5x^3+5x-2\) 
determinare il 
suo segno.
\end{esercizio}

\begin{esercizio}
\label{ese:4.34}
Dato il trinomio \(p(x)=9x^2+x^4-10\) stabilire se esiste almeno un numero 
naturale che lo renda negativo.
\end{esercizio}

\begin{esercizio}
\label{ese:4.35}
Nell'insieme dei valori reali che rendono positivo il trinomio 
\(p(x)=2x^5-12x^3-14x\) vi sono solo due numeri interi negativi?
\end{esercizio}

\begin{esercizio}
\label{ese:4.36}
\(x\in (-1;+\infty )\Rightarrow p(x)=x^5-2x^2-x+2>0\) Vero o falso?
\end{esercizio}

\begin{esercizio}
\label{ese:4.37}
Nell'insieme dei valori reali che rendono negativo \(p(x)=(2x-1)^3-(3-6x)^2\) 
appartiene un valore razionale che lo annulla. Vero o falso?
\end{esercizio}

\begin{esercizio}[*]
\label{ese:4.38}
Risolvi le seguenti disequazioni di grado superiore al secondo.
\begin{enumeratea}
\item \((1-x)(2-x)(3-x)>0\) \sol{x<1 \sor 2<x<3}
\item \((2x-1)(3x-2)(4x-3) \leqslant 0\) 
\sol{\dfrac 2 3 \leqslant x \leqslant \dfrac 3 4 \sor x \leqslant \dfrac 1 2}
\item \(-2x(x-1)(x+2)>0\) \sol{x<-2 \sor 0<x<1}
\item \( \left(x^4-4x^2-45\right)\cdot \left(4x^2-4x+1\right)>0 \) 
\sol{}
\item \(3x(x-2)(x+3)(2x-1) \leqslant 0\) 
\sol{\dfrac 1 2 \leqslant x \leqslant 2 \sor -3 \leqslant x \leqslant 0}
\item \(\left(x^2+1\right)(x-1)(x+2)>0\) \sol{x<-2 \sor x>1}
\item \(\left(1-9x^2\right)\left(9x^2-3x\right)2x>0\) 
\sol{x<-1/3}
\item \(\left(16x^2-1\right)\left(x^2-x-12\right)>0\) 
\sol{-\dfrac 1 4<x<\dfrac 1 4 \sor x<-3 \sor x>4}
\item \(-x\left(x^2-3x-10\right)\left(x^2-9x+18\right) \leqslant 0\)
\sol{3 \leqslant x \leqslant 5 \sor -2 \leqslant x \leqslant 0 \sor x \geqslant 
6}
\item \(x^2(x-1)\left(2x^2-x\right)\left(x^2-3x+3\right)>0\)
\sol{0<x<\dfrac 1 2 \sor x>1}
\item \(\left(x^2-1\right)\left(x^2-2\right)\left(x^2-3x\right)>0\)
\sol{x<-\sqrt 2 \sor 1<x<\sqrt 2 \sor -1<x<0 \sor x>3}
\item \(x^3-x^2+x-1>0\) \sol{x>1}
\item \(x^3-5x^2+6<0\) \sol{3-\sqrt 3<x<3+\sqrt 3 \sor x<-1}
\item \(\left(5x^3-2x^2\right)\left(3x^2-5x\right) \geqslant 0\) 
\sol{0 \leqslant x \leqslant \dfrac 2 5 \sor x \geqslant \dfrac 5 3}
\item \(x^4-2x^3-x+2>0\) \sol{x<1 \sor x>2}
\item \(x^4+x^2-9x^2-9 \leqslant 0\) \sol{-3 \leqslant x \leqslant 3}
\item \(25x^4-9>0\) 
\sol{x<-\dfrac{\sqrt{15}} 5 \sor x>\dfrac{\sqrt{15}} 5}
% \item \(x^3-1 \geqslant 2x(x-1)\) \sol{x \geqslant 1}
% \item \(x^4-1>x^2+1\) \sol{x<-\sqrt 2 \sor x>\sqrt 2}
\item \(\left(x^2+x\right)^2+2\left(x+1\right)^2 \geqslant 0\) 
\sol{\R}
% \item \((x+1)\left(x-\dfrac 1 2\right)(x+2)<0\) 
% \sol{-1<x<\dfrac 1 2 \sor x<-2}
% \item \(\left(x^2-4\right)(x-2) \geqslant 0\) 
% \sol{x \geqslant -2}
% \item \((x-7)\left(x^2-7x+10\right)<0\) 
% \sol{5<x<7 \sor x<2}
% \item \(\left(x^2-4\right)\left(x^2-9\right) \geqslant 0\) 
% \sol{x \leqslant -3 \sor -2 \leqslant x \leqslant 2 \sor x \geqslant 3}
\end{enumeratea}
\end{esercizio}
% 
% \begin{esercizio}[*]
% \label{ese:4.44}
% Risolvi le seguenti disequazioni di grado superiore al secondo.
% \begin{htmulticols}{2}
% \begin{enumeratea}
% \item \(\left(x^4+4x^3-12x^2\right)\left(x+3\right) \geqslant 0\)
% \item \((x-4)^3-(x-4)^2-2x+10>2\)
% \item \(x^3-1 \geqslant 0\)
% \item \(\left(x^4+4x^3-12x^2\right)\left(x+3\right) \geqslant 0\)
% \end{enumeratea}
% \end{htmulticols}
% \end{esercizio}
% 
% \begin{esercizio}[*]
% \label{ese:4.45}
% Risolvi le seguenti disequazioni di grado superiore al secondo.
% \begin{enumeratea}
% \item \((x+3)(x+4)(x+5)(5-x)(4-x)(3-x)>0\)
% \item \((x^2-2x)(x^2+1)>0\)
% \item \((8-2x^2)(3x-x^2+4)<0\)
% \item \((6x^2-6)(100x^2+100x)<0\)
% \end{enumeratea}
% \end{esercizio}
% 
% \paragraph{4.44.} a)~\(x=0 \sor -6 \leqslant x \leqslant -3 \sor x \geqslant 
% 2\),\quad 
% b)~\(3<x<4 \sor 
% x>6\),\quad c)~\(x \geqslant 1\),\protect\\
% d)~\(-9<x<-6 \sor -\dfrac 1 2<x<3\)
% 
% \paragraph{4.45.} a)~ \(-5<x<-4 \sor -3<x<3 \sor 4<x<5\),\quad b)~\(x<0 \sor 
% x>2\),\protect\\
% c)~\(-2<x<-1 \sor 2<x<4\),\quad d)~\(0<x<1\)
% 
% \begin{esercizio}[*]
% \label{ese:4.46}
% Risolvi le seguenti disequazioni di grado superiore al secondo.
% \begin{htmulticols}{2}
% \begin{enumeratea}
% \item \((1+x^2)(3x^2+x)<0\)
% \item \((x^2+3x+3)(4x^2+3)>0\)
% \item \((125+4x^2)(128+2x^2)<0\)
% \item \((x^2+4x+4)(x^2-4x+3)>0\)
% \end{enumeratea}
% \end{htmulticols}
% \end{esercizio}
% \newpage
% \begin{esercizio}[*]
% \label{ese:4.47}
% Risolvi le seguenti disequazioni di grado superiore al secondo.
% \begin{htmulticols}{2}
% \begin{enumeratea}
% \item \((x^2-5x+8)(x^2-2x+1)>0\)
% \item \((-2x+1)(3x-x^2)>0\)
% \item \((4x^2-3x)(x^2-2x-8)<0\)
% \item \((4x-x^2+5)(x^2-9x+20)<0\)
% \end{enumeratea}
% \end{htmulticols}
% \end{esercizio}
% 
% \begin{esercizio}[*]
% \label{ese:4.48}
% Risolvi le seguenti disequazioni di grado superiore al secondo.
% \begin{htmulticols}{2}
% \begin{enumeratea}
% \item \((5+2x)(-2x^2+14x+16)<0\)
% \item \((5x-2x^2-10)(x^2+3x-28)>0\)
% \item \((x^2-6x+9)(8x-7x^2)>0\)
% \item \((3x^2+2x-8)(6x^2+19x+15)<0\)
% \end{enumeratea}
% \end{htmulticols}
% \end{esercizio}
% 
% \begin{esercizio}[*]
% \label{ese:4.49}
% Risolvi le seguenti disequazioni di grado superiore al secondo.
% \begin{htmulticols}{2}
% \begin{enumeratea}
% \item \((3x^2-5x-2)(4x^2+8x-5)>0\)
% \item \((4x-4)(2x^2-3x+2)<0\)
% \item \((2x-4)(2x^2-3x-14)>0\)
% \item \((-7x+6)(x^2+10x+25)<0\)
% \end{enumeratea}
% \end{htmulticols}
% \end{esercizio}
% 
% \begin{esercizio}[*]
% \label{ese:4.50}
% Risolvi le seguenti disequazioni di grado superiore al secondo.
% \begin{htmulticols}{2}
% \begin{enumeratea}
% \item \((-3+3x)(x^3-4x^2)>0\)
% \item \(\left(x^2+1\right)\left(x^2-1\right)>0\)
% \item \((1-x)(2-x)^2 \leqslant 0\)
% \item \(-x\left(x^2+1\right)(x+1) \geqslant 0\)
% \item \((x+1)^2\left(x^2-1\right)<0\)
% \item \((x^2-4)(2x-50x^2) \geqslant 0\)
% \end{enumeratea}
% \end{htmulticols}
% \end{esercizio}
% 
% \begin{esercizio}[*]
% \label{ese:4.51}
% Risolvi le seguenti disequazioni di grado superiore al secondo.
% \begin{htmulticols}{2}
% \begin{enumeratea}
% \item \((x-4)(2x^2+x-1) \geqslant 0\)
% \item \(-3x^3+27>0\)
% \item \(3x^3+27>0\)
% \item \(x^3+3x^2+3x+1 \leqslant 0\)
% \item \(x^3-6x+9<0\)
% \item \(x^5+1>x\left(x^3+1\right)>0\)
% \end{enumeratea}
% \end{htmulticols}
% \end{esercizio}
% 
% \paragraph{4.46.} a)~\(-\dfrac 1 3<x<0\),\quad b)~\(\IS=\R\),\quad 
% c)~\(\IS=\emptyset \),\quad d)~\(x<-2 \sor -2<x<1 \sor x>3\)
% 
% \paragraph{4.47.} a)~\(x<1 \sor x>1\),\quad b)~
% \(0<x<\dfrac 1 2 \sor x>3\),\quad 
% c)~\(-2<x<0 \sor \dfrac 3 4<x<4\),\protect\\
% d)~\(x<-1 \sor 4<x<5 \sor x>5\)
% 
% \paragraph{4.48.} a)~\(-\dfrac 5 2<x<-1 \sor x>8\),\quad b)~\(-7<x<4\),\; 
% c)~\(0<x<\dfrac 8 7\),\; d)~
% \(-2<x<-\dfrac 5 3 \sor -\dfrac 3 2<x<\dfrac 4 3\)
% 
% \paragraph{4.49.} a)~
% \(x<-\dfrac 5 2 \sor -\dfrac 1 3<x<\dfrac 1 2 \sor x>2\),\quad 
% b)~\(x<1\),\quad c)~\(-2<x<2 \sor x>\dfrac 7 2\),\quad d)~\(x>6/7\)
% 
% \paragraph{4.50.} a)~\(\IS=x\in \R| x<0 \sor 0<x<1 \sor x>4\)
% 
% \paragraph{4.51.} d)~\(x \leqslant -1\),\quad e)~\(x<-3\)

% \begin{esercizio}
% \label{ese:4.52}
% Risolvi le seguenti disequazioni di grado superiore al secondo.
% \begin{htmulticols}{2}
% \begin{enumeratea}
% \item \(x^3-7x^2+4x+12 \geqslant 0\)
% \item \(x^3+5x^2-2x-24<0\)
% \item \(6x^3+23x^2+11x-12 \leqslant 0\)
% \item \(4x^3+4x^2-4x-4 \geqslant 0\)
% \item \(-6x^3-30x^2+192x-216<0\)
% \item \(81x^4-1 \leqslant 0\)
% \end{enumeratea}
% \end{htmulticols}
% \end{esercizio}
% 
% \begin{esercizio}
% \label{ese:4.53}
% Risolvi le seguenti disequazioni di grado superiore al secondo.
% \begin{htmulticols}{2}
% \begin{enumeratea}
% \item \(3x^5+96<0\)
% \item \(x^4-13x^2+36<0\)
% \item \(9x^4-37x^2+4 \geqslant 0\)
% \item \(-4x^4+65x^2-16<0\)
% \item \(x^6-4x^3+3 \geqslant 0\)
% \item \(x^8-x^4-2<0\)
% \end{enumeratea}
% \end{htmulticols}
% \end{esercizio}
% 
% \begin{esercizio}
% \label{ese:4.54}
% Risolvi le seguenti disequazioni di grado superiore al secondo.
% \begin{enumeratea}
% \item \(\dfrac 2 3x^3>\dfrac 9 4\)
% \item \( (2x-1)^2 \geqslant x^2\left(4x^2-4x+1\right) \)
% \item \( (x+1)\left(x^2-1\right)>\left(x^2-x\right)(x-1)^2 \)
% \item \( -4x\left(x^2+7x+12\right)\left(x^2-25\right)(4-x)>0 \)
% \item \( (x-5x^2)(x^4-3x^3+5x^2) \geqslant 0 \)
% \item \( (4+7x^2)\left[x^2-(\sqrt 2+\sqrt 3)x+\sqrt 6\right]<0 \)
% \end{enumeratea}
% \end{esercizio}
% \newpage
% \begin{esercizio}
% \label{ese:4.55}
% Risolvi le seguenti disequazioni di grado superiore al secondo.
% \begin{htmulticols}{2}
% \begin{enumeratea}
% \item \( (x^3-9x)(x-x^2)(4x-4-x^2)>0 \)
% \item \( x\left|x+1\right|\cdot (x^2-2x+1) \geqslant 0 \)
% \item \(16x^4-1 \geqslant 0\)
% \item \(16x^4+1 \leqslant 0\)
% \item \(-16x^4-1>0\)
% \item \(-16x^4+1>0\)
% \end{enumeratea}
% \end{htmulticols}
% \end{esercizio}
% 
% \begin{esercizio}
% \label{ese:4.56}
% Risolvi le seguenti disequazioni di grado superiore al secondo.
% \begin{htmulticols}{3}
% \begin{enumeratea}
% \item \(1-16x^4<0\)
% \item \(27x^3-8 \geqslant 0\)
% \item \(8x^3+27<0\)
% \item \(4x^4+1 \geqslant 0\)
% \item \(4x^4-1 \geqslant 0\)
% \item \(1000x^3+27>0\)
% \end{enumeratea}
% \end{htmulticols}
% \end{esercizio}
% 
% \begin{esercizio}
% \label{ese:4.57}
% Risolvi le seguenti disequazioni di grado superiore al secondo.
% \begin{htmulticols}{3}
% \begin{enumeratea}
% \item \(10000x^4-1 \geqslant 0\)
% \item \(x^7+7<0\)
% \item \(x^3-8 \geqslant 0\)
% \item \(9x^4-4 \geqslant 0\)
% \item \(x^6+\sqrt 6 \leqslant 0\)
% \item \(0,1x^4-1000 \geqslant 0\)
% \item \(x^4-9 \geqslant 0\)
% \item \(x^4+9 \leqslant 0\)
% \item \(-x^4+9 \leqslant 0\)
% \end{enumeratea}
% \end{htmulticols}
% \end{esercizio}

% \subsection*{4.5 - Disequazioni fratte}
\numnameref{sec:diseq_fratte}

\begin{esercizio}[*]
\label{ese:4.58}
Determinare l'Insieme Soluzione delle seguenti disequazioni fratte.
\begin{htmulticols}{2}
\begin{enumeratea}
\item \(\dfrac{x+2}{x-1}>0\) \sol{x<2 \sor x>1}
% \item \(\dfrac{x+3}{4-x}>0\) \sol{-3<x<4}
% \item \(\dfrac{x+5}{x-7}>0\) \sol{x<-5 \sor x>7}
\item \(\dfrac{2-4x}{3x+1} \geqslant 0\) \sol{-\dfrac 1 3<x \leqslant \dfrac 1 
2}
\item \(\dfrac{x^2-4x+3}{4-7x} \geqslant 0\) 
\sol{x<\dfrac 4 7 \sor 1 \leqslant x \leqslant 3}
\item \(\dfrac{x+5}{x^2-25}>0\) \sol{x>5}
\item \(\dfrac{x^2-1}{x-2}>0\) \sol{-1<x<1 \sor x>2}
\item \(\dfrac{x^2-4x+3}{x+5}<0\) \sol{x<-5 \sor 1<x<3}
% \item \(\dfrac{-x^2+4x-3}{x+5}>0\) \sol{x<-5 \sor 1<x<3}
\item \(\dfrac{x^2+1}{x^2-2x}>0\) \sol{x<0 \sor x>2}
\item \(\dfrac{9-x^2}{2x^2-x-15}>0\) \sol{-3<x<-\dfrac 5 2}
\item \(\dfrac{x^2-7x}{-x^2-8}>0\) \sol{0<x<7}
\item \(\dfrac{x+2}{x-1} \leqslant 0\) \sol{1< x \leqslant 2}
\item \(\dfrac 1{x^2+2x+1}>0\) \sol{\R-\{-1\}\}}
\item \(\dfrac{-3}{-x^2-4x-8}>0\) \sol{\R}
\item \(\dfrac{x^2+2x+3}{-x^2-4}>0\) \sol{\emptyset}
% \item \(\dfrac{3x-12}{x^2-9}>0\) \sol{-3<x<3 \sor x>4}
\item \(\dfrac{5-x}{x^2-4}>0\) \sol{x<-2 \sor 2<x<5}
\end{enumeratea}
\end{htmulticols}
\end{esercizio}

\begin{esercizio}[*]
\label{ese:4.61}
Determinare l'Insieme Soluzione delle seguenti disequazioni fratte.
\begin{htmulticols}{2}
\begin{enumeratea}
\item \(\dfrac{3x-x^2-2}{2x^2+5x+3}>0\) \sol{-\dfrac 3 2<x<-1 \sor 
1<x<2}
\item \(\dfrac{4-2x}{x^2-2x-8}>0\) \sol{x<-2 \sor 2<x<4}
\item \(\dfrac{x^2-4x+3}{5-10x}>0\) \sol{x<\dfrac 1 2 \sor 
1<x<3}
\item \(\dfrac{x^2+3x+10}{4-x^2}>0\) \sol{-2<x<2}
\item \(\dfrac{x^2-3x+2}{4x-x^2-5}>0\) \sol{1<x<2}
\item \(\dfrac{x^2+2}{25-x^2}>0\) \sol{-5<x<5}
\item \(\dfrac{3x^2-2x-1}{4-2x}>0\) \sol{x<-\dfrac 1 3 \sor 
1<x<2}
\item \(\dfrac{x+2}{x^2+4x+4}>0\) \sol{x>-2}
\item \(\dfrac{x+2}{x^2+4x+2}>0\) \sol{x<1 \sor 3<x<5}
\item \(\dfrac{-x^2+2x+8}{-x-1}<0\) \sol{x<-2 \sor -1<x<4}
\item \(\dfrac{x^2+3x+2}{25-x^2}>0\) \sol{-5<x<-2 \sor 
-1<x<5}
\item \(\dfrac{x^2+4x+3}{3x-6}>0\) \sol{-3<x<-1 \sor x>2}
\item \(\dfrac{5-x}{x^2-4x+3}>0\) \sol{x<-\dfrac 3 4 \sor 
1<x<4}
\item \(\dfrac{1-x^2}{x^2+2x+3}<0\) \sol{x<-1 \sor x>1}
% \item \(\dfrac{x^2-9}{x^2-5x}>0\) \sol{x<-3 \sor 0<x<3 \sor 
% x>5}
% \item \(\dfrac{x^2-x-2}{x-x^2+6}>0\) \sol{-2<x<-1 \sor 
% 2<x<3}
% \item \(\dfrac{x^2-5x+6}{-3x+7}<0\) \sol{2<x<\dfrac 7 3 \sor 
% x>3}
% \item \(\dfrac{2x+8}{x^2+4x-12}>0\) \sol{-6<x<-4 \sor x>2}
\end{enumeratea}
\end{htmulticols}
\end{esercizio}

% \pagebreak %--------------------------------------------------

\begin{esercizio}[*]
\label{ese:4.64}
Determinare l'Insieme Soluzione delle seguenti disequazioni fratte.
\begin{enumeratea}
% \item \(\dfrac{x^2-2x-63}{4x+5-x^2}>0\) \sol{-7<x<-1 \sor 
% 5<x<9}
% \item \(\dfrac{4-x^2+3x}{x^2-x}>0\) \sol{-1<x<0 \sor 1<x<4}
% \item \(\dfrac{x^2-2x}{5-x^2}>0\) 
% \sol{-\sqrt 5<x<0 \sor 2<x<\sqrt 5}
\item \(\dfrac{x^2-x-2}{-3x^2+3x+18} \leqslant 0\) 
\sol{x<-2 \sor -1 \leqslant x \leqslant 2 \sor x>3}
\item \(\dfrac{x^2-8x+15}{x^2+3x+2}>0\) 
\sol{x<-2 \sor -1<x<3 \sor x>5}
\item \(\dfrac{4x+7}{3x^2-x-2}>0\) 
\sol{-\dfrac 7 4<x<-\dfrac 2 3 \sor x>1}
\item \(\dfrac{-x^2-4x-3}{6x-x^2}>0\) 
\sol{x<-3 \sor -1<x<0 \sor x>6}
\item \(\dfrac{5x+x^2+4}{6x^2-6x}>0\) 
\sol{x<-4 \sor -1<x<0 \sor x>1}
\item \(\dfrac{9-x^2}{x^2+5x+6}\cdot \dfrac{6x-2x^2}{4-x^2}>0\) 
\sol{0<x<2}
\item \(\dfrac{2x-4x^2}{x^2+x-12}\cdot \dfrac{16-x^2}{5x-x^2} \leqslant 0\) 
\sol{x \leqslant \dfrac 1 2 \sor 3<x \leqslant 4 \sor x>5 \wedge 
x\neq 0 \wedge x\neq -4}
\item \(\dfrac{1-x^2}{x^2} \leqslant \dfrac 1{x^2}-x^2-\dfrac 1 2\) 
\sol{-\dfrac{\sqrt 2} 2 \leqslant x \leqslant \dfrac{\sqrt 2} 2 \wedge 
x\neq 0}
\item \(\dfrac{x+2}{x-1} \geqslant \dfrac{24}{x+1}-\dfrac x{x^2-1}\) 
\sol{x<-1 \sor 1<x \leqslant 10-\sqrt{74} \sor x \geqslant 10+\sqrt{74}}
\item \(\dfrac 1 x+\dfrac 1{x-1}+\dfrac 1{x+1}<\dfrac{2x+1}{x^2-1}\) 
\sol{-\dfrac{\sqrt 2} 2 \leqslant x \leqslant \dfrac{\sqrt 2} 2}
\item \(\dfrac x{x+2} \geqslant \dfrac{x-4}{x^2-4}\) 
\sol{x<-2 \sor x>2}
\item \(\dfrac{4x+1}{x^2-9}+\dfrac{1-x}{x+3}<6-\dfrac x{x-3}\) 
\sol{x<-3 \sor -\dfrac{13} 6<x<3 \sor x>4}
\item \(\dfrac{x+1}{2x-1}+\dfrac 3{4x+10} \geqslant 1-\dfrac{2x+2}{4x^2+8x-5}\) 
\sol{-\dfrac 5 2<x \leqslant -\dfrac 3 2 \sor \dfrac 1 2<x \leqslant \dfrac 7 
2}
\item \(\dfrac{2x+5}{(2x+4)^2} \geqslant \dfrac 2{2x+4}\) 
\sol{x \leqslant -\dfrac 3 2 \wedge x \neq -2}
\item \(\dfrac{10x^2}{x^2+x-6}+\dfrac x{2-x}-1 \leqslant \dfrac 5{x+3}\) 
\sol{-3<x<2}
\item \(\dfrac{5x+20}{5x+5}+\dfrac{2x-8}{2x-2} \geqslant 2\) 
\sol{-1<x<1}
\item \(\dfrac 8{8x^2-8x-70}-\dfrac 4{4x^2-4x-35}>\dfrac{8x+8}{4x^2-20x+21}\)
\sol{\dfrac 3 2<x<\dfrac 7 2 \sor x<-1}
\item \(\dfrac{4x^2-8x+19}{8x^2-36x+28}-\dfrac{2x-5}{4x-4} \geqslant 
\dfrac{8x+12}{8x-28}\)
\sol{\dfrac 1 2 \leqslant x<\dfrac 7 2 \wedge x \neq 1}
\end{enumeratea}
\end{esercizio}

\begin{esercizio}[*]
\label{ese:4.68}
Assegnate le due funzioni \(f_1=\dfrac{x^2+1}{2x-x^2}\) e \(f_2=\dfrac 1 
x+\dfrac 
1{x-2}\) stabilire per quali valori della variabile indipendente si ha 
\(f_1 \geqslant 
f_2\)
\sol{-1-\sqrt 2 \leqslant x<0 \sor -1+\sqrt 2 \leqslant x<2}
\end{esercizio}

\begin{esercizio}
\label{ese:4.69}
Spiegare perché l'espressione letterale 
\(E=\dfrac{1-\dfrac{x^2}{x^2-1}}{2+\dfrac{3x-1}{1-x}}\) è sempre positiva nel 
suo 
dominio.
\end{esercizio}

\begin{esercizio}
\label{ese:4.70}
Per quali valori di x la funzione \(y=\dfrac{(x-1) \cdot x-2}{5x^2-x-4}\) 
è maggiore o uguale a 1.
\sol{-\dfrac 3 2 \leqslant x<-\dfrac 4 5}
\end{esercizio}

\begin{esercizio}[*]
\label{ese:4.71}
\( x \), \( x+2 \), \( x+4 \) sono tre numeri naturali. Determinate in \( 
\N \) il 
più piccolo numero che rende vera la proposizione: ``il doppio del primo 
aumentato del prodotto degli altri due è maggiore della differenza tra il 
doppio 
del terzo e il quadrato del secondo''
\sol{5}
\end{esercizio}

\begin{esercizio}
\label{ese:4.72}
Date chiare e sintetiche motivazioni alla verità della seguente proposizione: 
``il segno della frazione \(f=\dfrac{9-x^2+3x}{2+x^2}\) non è mai positivo e 
la 
frazione non ha zeri reali''.
\end{esercizio}

\begin{esercizio}
\label{ese:4.73}
Stabilire se basta la condizione \(x\neq 1\wedge x\neq -1\) per rendere 
positiva 
la frazione \(~f=~\dfrac{x^3-1}{x^4-2x^2+1}\)
\end{esercizio}

\begin{esercizio}
\label{ese:4.74}
Determinare per quali valori reali la frazione 
\(f=\dfrac{(x+1)^2}{4x^2-12x+9}\) 
risulta non superiore a 1.
\end{esercizio}

% \subsection*{4.6 - Sistemi di disequazioni}
\numnameref{sec:diseq_sistemi}

\begin{esercizio}
\label{ese:21.33}
Sulla retta reale rappresenta l'insieme soluzione~\(S_{1}\)
dell'equazione:
\[\dfrac{1}{6}+\dfrac{1}{4}\cdot (5x+3)=2+\dfrac{2}{3}\cdot (x+1)\]

e l'insieme soluzione~\(S_{2}\) della disequazione:
\[\dfrac{1}{2}-2\cdot\left(\dfrac{1-x}{4}\right)\ge~3-\dfrac{6-2x}{3}-\dfrac{x
}{2}.\]

È vero che~\(S_{1}\subset S_{2}\)?
\end{esercizio}

\begin{esercizio}[*]
\label{ese:21.34}
Determina i numeri reali che verificano il sistema:
\(\sistema{x^{2}\le~0 \\ 2-3x\ge~0}\)
\sol{x = 0}
\end{esercizio}

\begin{esercizio}
\label{ese:21.35}
L'insieme soluzione del sistema:
\(\sistema{
(x+3)^{3}-(x+3)\cdot (9x-2)>x^{3}+27\\
\dfrac{x+5}{3}+3+\dfrac{2\cdot (x-1)}{3}<x+1
}\) è:
\begin{htmulticols}{2}
\fbox{A}\quad~\(\left\{x\in \R/x>3\right\}\)

\fbox{B}\quad~\(\left\{x\in \R/x>-3\right\}\)

\fbox{C}\quad~\(\left\{x\in \R/x<-3\right\}\)

\fbox{D}\quad~\(\IS=\emptyset \)

\fbox{E}\quad~\(\left\{x\in\R/x<3\right\}\)
\end{htmulticols}

\end{esercizio}

\begin{esercizio}
\label{ese:21.36}
Attribuire il valore di verità alle seguenti proposizioni:

\begin{enumeratea}
\item il quadrato di un numero reale è sempre positivo;
\item l'insieme complementare di~\(A=\{x\in\R/x>-8\}\text{ è 
}B=\{x\in\R/x<-8\}\)
\item il monomio~\(-6x^{3}y^{2}\) assume valore positivo per tutte le coppie 
dell'insieme~\(\R^{+}\times\R^{+}\)
\item nell'insieme~\(\Z\) degli interi relativi il 
sistema~\(\sistema{x+1>0\\8x<0}\) non ha 
soluzione;
\item l'intervallo~\(\left[-1,\left.-{\dfrac{1}{2}}\right)\right.\) 
rappresenta l'\(\IS\) del sistema~\(\sistema{1+2x<0 
\\\dfrac{x+3}{2} \leqslant x+1}\)
\end{enumeratea}
\end{esercizio}

% \newpage %-----------------------------------------------

\begin{esercizio}[*]
\label{ese:21.37}
Risolvi i seguenti sistemi di disequazioni.
\begin{htmulticols}{2}
\begin{enumeratea}
\item \(\sistema{
3-x>x\\
2x>3
  }\)
\sol{\emptyset}
\item \(\sistema{
3x\le~4\\
5x \geqslant -4
}\)
\sol{-{\dfrac{4}{5}} \leqslant x\le\dfrac{4}{3}}
\item \(\sistema{
2x>3\\
3x\le~4
  }\)
\sol{\emptyset}
\item \(\sistema{
3x-5<2\\
x+7<-2x
}\)
\sol{x<-{\dfrac{7}{3}}}
\item \(\sistema{
3-x \geqslant x-3\\
-x+3\ge~0
}\)
\sol{x\le~3}
\item \(\sistema{
-x-3\le~3\\
3+2x\ge~3x+2
}\)
\sol{-6 \leqslant x\le~1}
\item \(\sistema{
2x-1>2x \\
3x+3\le~3
}\)
\sol{\emptyset}
\item \(\sistema{
2x+2<2x+3\\
2(x+3)>2x+5
}\)
\sol{\R}
\item \(\sistema{
-3x>0\\
-3x+5\ge~0\\
-3x\ge-2x
}\)
\sol{x<0}
\item {\arraystretch{2}\(\sistema{
-{\dfrac{4}{3}}x\ge\dfrac{2}{3}\\
-{\dfrac{2}{3}}x\le\dfrac{1}{9}
}\)}
\sol{\emptyset}
\item \(\sistema{
3+2x>3x+2 \\
5x-4\le~6x-4\\
-3x+2 \geqslant -x-8
}\)
\sol{0 \leqslant x<1}
\item \(\sistema{
4x+4\ge~3\cdot\left(x+\dfrac{4}{3}\right)\\
4x+4\ge~2\cdot (2x+2)
}\)
\sol{x\ge~0}
\item \(\sistema{
3(x-1)<2(x+1)\\
x-\dfrac{1}{2}+\dfrac{x+1}{2}>0
}\)
\sol{0<x<5}
\end{enumeratea}
\end{htmulticols}
\end{esercizio}

\begin{esercizio}[*]
\label{ese:4.75}
Risolvere i seguenti sistemi di disequazioni.
\begin{htmulticols}{2}
\begin{enumeratea}
\item \(\sistema{x^2-4>0\\x-5 \leqslant 0}\) 
\sol{x<-2 \sor 2<x}
\item \(\sistema{x^2-4x+3 \leqslant 
0\\x-2x^2<-10}\)
\sol{\dfrac 5 2<x \leqslant 3}
\item \(\sistema{4x-x^2>0\\3x^2(x-3)>0}\)
\sol{3<x<4}
\item \(\sistema{x^2+5x+6 \leqslant 0\\2x+5 \leqslant 
0}\)
\sol{-3 \leqslant x \leqslant -\dfrac 5 2}
\item \(\sistema{3x-x^2-2 \leqslant 0\\x^2>49}\)
\sol{x<-7 \sor x>7}
\item \(\sistema{3x-2>0\\x^2-1>0\\2x-x^2<0}\)
\sol{x>2}
\item \(\sistema{x^2-4x+4 \geqslant 0\\x<6}\)
\sol{x<6}
\item \(\sistema{x^2+6x+9<0\\x<2\\x^2+1>0}\)
\sol{\emptyset}
\item \(\sistema{x^2+6x+9 \leqslant 0\\x<2}\)
\sol{x=-3}
\item \(\sistema{4x-x^2-3<0\\3x \geqslant 2}\)
\sol{\dfrac 2 3 \leqslant x<1 \sor x>3}
\item \(\sistema{2x^2<8\\-x^2+5x>-6\\x^2(9-x^2) \leqslant 0}\)
\sol{x=0}
\item \(\sistema{(x^2-4x+3)(2x-4)>0\\2x-x^2 \leqslant 1}\)
\sol{1<x<2 \sor x>3}
\item \(\sistema{(3-x)(x^2-4)(x^2-2x-8)<0\\x^2-64 \leqslant 0}\)
\sol{2<x<3 \sor 4<x \leqslant 8}
\item \(\sistema{2x^2-x-1 \leqslant 0\\3x+7>0\\x^2-10x+9 \leqslant 0}\)
\sol{x=1}
\item \(\sistema{2x^2-x-1<0\\3x+7>0\\x^2-10x+9 \leqslant }\)
\sol{\emptyset}
\item \(\sistema{x^2-10x+25>0\\x<7}\)
\sol{x<5 \sor 5<x<7}
\item \(\sistema{x^2-10x+25 \geqslant 0\\x<7}\)
\sol{x<7}
% \item \(\sistema{x^4-8 \geqslant 1 \\ \dfrac{5-x} x<\dfrac 1}\)
% \sol{x \leqslant -\sqrt 3}
\item \(\sistema{x^2-4x+3 \leqslant 0\\x^2-4>0 \\x^2+1>0 \\x-1>0}\)
\sol{2<x \leqslant 3}
\item \(\sistema{x^2-5x+6 \leqslant 0\\x^2-1>0 \\x^2+1<0 \\x-1>0}\)
\sol{\emptyset}
\item \(\sistema{
  x^2-2x+1 \geqslant 0 \\
  x^2+5x \geqslant 0 \\
  x^2+1>0 \\
  x^2-2x+7>0}\)
\sol{x \leqslant -5 \sor x \geqslant 0}
% \item 
% \(\sistema{x^2-3x+2>0\\x^2-3x+2<0\\2x^2-x-1>0\\x^2-2x>0}\)
% \sol{\emptyset}
% \item \(\sistema{
%   x^2-3x+2 \leqslant 0\\
%   x^2-4x+4 \leqslant 0\\
%   x^2-3x+2 \geqslant 0\\
%   x^2-4x+4 \geqslant 0}\)
% \sol{x=2}
\end{enumeratea}
\end{htmulticols}
\end{esercizio}

% \newpage %--------------------------------------------

\begin{esercizio}[*]
\label{ese:21.41}
Risolvi i seguenti sistemi di disequazioni.
\begin{enumeratea}
\item \(\sistema{x^2-4x+4>0\\x \leqslant 6\\1-x^2 \leqslant 
0}\)
\sol{x \leqslant -1 \sor 1 \leqslant x<2 \sor 2<x \leqslant 6}
\item {\(\sistema{
\dfrac{2x+3}{3}>x-1\\[3mm]
\dfrac{x-4}{5}<\dfrac{2x+1}{3}}\)}
\sol{-{\dfrac{17}{7}}<x<6}
\item {\(\sistema{
16(x+1)-2+(x-3)^{2} \leqslant (x+5)^{2}\\[3mm]
\dfrac{x+5}{3}+3+2\cdot\dfrac{x-1}{3} \leqslant x+4
}\)}
\sol{\R}
\item \(\sistema{
x+\dfrac{1}{2}<\dfrac{1}{3}(x+3)-1\\
(x+3)^{2} \geqslant (x-2)(x+2)
}\)
\sol{-{\dfrac{13}{6}} \leqslant x<-{\dfrac{3}{4}}}
\item {\(\sistema{
2\left(x-\dfrac{1}{3}\right)+x>3x-2\\[3mm]
\dfrac{x}{3}-\dfrac{1}{2} \geqslant \dfrac{x}{4}-\dfrac{x}{6}
}\)}
\sol{x\ge~2}
\item \(\sistema{
\dfrac{3}{2}x+\dfrac{1}{4}<5\cdot\left(\dfrac{2}{3}x-\dfrac{1}{2}\right)\\
x^2-2x+1\ge~0
}\)
\sol{x>\dfrac{3}{2}}
\item {\(\sistema{
3\left(x-\dfrac{4}{3}\right)+\dfrac{2-x}{3}+x-\dfrac{x-1}{3}>0\\[3mm]
\left[1-\dfrac{1}{6}(2x+1)\right]+\left(x-\dfrac{1}{2}\right)^{2}<(x+1)^{2}
+\dfrac{1}{3}(1+2x)
 }\)}
\sol{x>\dfrac{9}{10}}
% \item {\(\sistema{
% \left(x-\dfrac{1}{2}\right)\left(x+\dfrac{1}{2}\right)>\left(x-\dfrac{1}{2}
% \right)^{2}\\[3mm]
% \left(x-\dfrac{1}{2}\right)\left(x+\dfrac{1}{2}\right)<\left(x-\dfrac{1}{2}
% \right)^{2}+\left(x+\dfrac{1}{2}\right)^{2}
% }\)}
% \sol{x>\dfrac{1}{2}}
\item \(\sistema{
   x^2-3x+2 \leqslant 0\\
   x^2-4x+4 \leqslant 0\\
   x^2-x+10>0\\
   x^2-2x \leqslant 0 }\)
\sol{x=2}
\item \(\sistema{
   \dfrac{4-x^2+3x}{x^2-x}>0 \\[3mm]
   \dfrac{x^2-x-2}{-3x^2+3x+18} \leqslant 0 }\)
\sol{3<x<4 \sor -1<x<0 \sor 1<x \leqslant 2}
\item \(\sistema{
   x^3-5x^2-14x \geqslant 0 \\ 
   \dfrac{2x+1}{2x} > \dfrac 3{x+1}}\)
\sol{2 \leqslant x<-1 \sor x \geqslant 7}
% \item \(\sistema{
%    \dfrac 1 x>\dfrac 1{x-3}\\
%    3x-1-2x^2<0\\
%    \dfrac{x^2-6x+5}{2-x}>0 }\)
% \sol{0<x<\dfrac 1 2 \sor 2<x<3}
% \item \(\sistema{x^2-2x+1>0\\
%        x^2+5x \geqslant 0 \\
%        x^2+x+23>0\\
%        x^2-2x+7>0
%    }\)
% \sol{x \leqslant -5 \sor 0 \leqslant x<1 \sor x>1}
\end{enumeratea}
\end{esercizio}

% \begin{esercizio}[*]
% \label{ese:4.82}
% Dato il sistema \(\sistema{{x(x-3)>3\left(\dfrac{x^2} 
% 2-2x\right)}\\{2+x\cdot \dfrac{3x-7} 3 \geqslant 5-\dfrac 1 
% 3x}}\) 
% determina i numeri naturali che lo risolvono. 
% \sol{3, 4, 5}
% \end{esercizio}

\begin{esercizio}[*]
\label{ese:4.83}
Per quali valori di \( x \) le due funzioni \(f_1=x^4-x^3+x-1\) e 
\(f_2=x^4-8x\) 
assumono contemporaneamente valore positivo? \sol{x<-1 \sor x>2}
\end{esercizio}
